\documentclass[11pt,a4paper]{article}

% -------------------------------------------------
% Encoding and language
% -------------------------------------------------
\usepackage[T1]{fontenc}
\usepackage[utf8]{inputenc}
\usepackage[magyar]{babel}
\usepackage{lmodern}
\usepackage{microtype}
\usepackage{setspace}
\usepackage{xcolor}
% -------------------------------------------------
% Math packages
% -------------------------------------------------
\usepackage{amsmath,amssymb,amsthm,mathtools}
\usepackage{mathrsfs}
\usepackage{stmaryrd}

% -------------------------------------------------
% Layout
% -------------------------------------------------
\usepackage[a4paper,margin=2.8cm]{geometry}
\usepackage{setspace}
\setstretch{1.08}

% -------------------------------------------------
% Hyperlinks
% -------------------------------------------------
\usepackage[hidelinks]{hyperref}

% -------------------------------------------------
% Enumerations
% -------------------------------------------------
\usepackage{enumitem}

% -------------------------------------------------
% Custom commands
% -------------------------------------------------
\newcommand{\Z}{\mathbb{Z}}
\newcommand{\Q}{\mathbb{Q}}
\newcommand{\R}{\mathbb{R}}
\newcommand{\C}{\mathbb{C}}
\newcommand{\N}{\mathbb{N}}
\newcommand{\Frac}{\operatorname{Frac}}
\newcommand{\Spec}{\operatorname{Spec}}
\newcommand{\length}{\operatorname{length}}
\newcommand{\ord}{\operatorname{ord}}
\newcommand{\m}{\mathfrak{m}}
\newcommand{\Ov}{\mathcal{O}_v}

% -------------------------------------------------
% Theorem styles
% -------------------------------------------------
\theoremstyle{definition}
\newtheorem{definition}{Definíció}[section]
\newtheorem{example}[definition]{Példa}
\newtheorem{remark}[definition]{Megjegyzés}

\theoremstyle{plain}
\newtheorem{theorem}[definition]{Tétel}
\newtheorem{proposition}[definition]{Állítás}
\newtheorem{corollary}[definition]{Következmény}

% -------------------------------------------------
% Title
% -------------------------------------------------
\title{%
Értékelések és értékelésgyűrűk\\
\large Különös tekintettel a diszkrét értékelésgyűrűkre
}

\author{Előadásjegyzet az Algebrai Görbék c. tárgyhoz}
\date{}

% -------------------------------------------------
% Document
% -------------------------------------------------
\begin{document}

\maketitle

\tableofcontents

\bigskip

% =================================================
{\section*{Előszó}
Az algebrai görbék tanulmányozása a matematika egy fontos területe, hisz hidat képez az komplex analízis, az algebra és a geometria között. Egy algebrai görbe nem más, mint egy $K$ test feletti polinom nulla halmaza, azaz egy $C = \{(x,y) \in K \times K \mid f(x,y) = 0\}$ halmaz, ahol $f$ egy
kétváltozós polinom, $K$-beli együtthatókkal.
\newline\indent Ezen jegyzet az algebrai görbék elméletének önálló, tisztán algebrai ismertetését adja. Ezt a megközelítést R. Dedekind,
L. Kronecker és H. M. Weber kezdeményezte a tizenkilencedik században majd E. Artin, H. Hasse, F. K. Schmidt és A. Weil továbbfejlesztették 
a huszadik század első felében. Ez a megközelítés csupán az algebrai testbővítések alapvető ismereteit feltételezi, beleértve a Galois-elméletet is. Egy másik előnye, hogy
az elmélet néhány főbb eredménye, mint például a Riemann—Roch-tétel, amely nem mellékesen ezen kurzus végcélja is, viszonylag hamar levezethető. 
\newline\indent Az első témakörünk az értékelések és az értékelésgyűrűk, melyekkel pontosan vizsgálhatjuk az algebrai görbék lokális viselkedését, lehetővé téve a zérus helyek, pólusok és szingularitások szisztematikus nyomon követését. A $C$ görbe egy $p$ pontja akkor sima pont, ha a $p$-ben lévő lokális gyűrű diszkrét értékelésgyűrű. Ellenkező esetben
$p$ szinguláris pont. %A $C$ görbe akkor sima, ha minden pontja sima.
Egy nem szinguláris algebrai görbe pontjai tehát megfelelnek a függvénytér diszkrét értékelésgyűrűinek. Ez a megfelelés pedig később lehetővé teszi például a divizorok definiálását és a Riemann–Roch-tétel megfogalmazását.


\section{Értékelések}

\subsection{Rendezett abel-csoportok}

\begin{definition}
Egy $\Gamma$ abel-csoportot \textit{(teljesen) rendezett abel-csoport}nak nevezünk,
ha el van látva egy teljes {$\leq$} rendezéssel, és a rendezés kompatibilis
a csoportművelettel:  ha $\alpha \leq \beta$, ahol $\alpha, \beta \in \Gamma$,  akkor minden $\gamma \in \Gamma$ esetén
\begin{equation*}
    \alpha + \gamma \leq \beta + \gamma.
\end{equation*}

Az értékelések értékkészletét egy ilyen $\Gamma$ csoport adja. A legfontosabb
eset számunkra a $\Gamma = \mathbb{Z}$ ellátva a szokásos rendezéssel. Ekkor beszélünk
diszkrét értékelésről.

Formálisan hasznos a $\Gamma$ csoportot kiegészíteni egy $\infty$ elemmel, amelyre
\begin{equation*}
    \alpha < \infty, \alpha + \infty = \infty + \alpha = \infty
\end{equation*}
minden $\alpha \in \Gamma $ esetén. Így kényelmesen definiálhatjuk $v(0) = \infty$  értéket.
\end{definition}

\subsection{Értékelések definíciója}

\begin{definition} \label{valuation}
Legyen $K$ test, $\Gamma$ rendezett abel-csoport.
Egy \textit{értékelés} a $K$ testen egy olyan
\[
v : K^\times \longrightarrow  \Gamma
\]
leképezés,  amelyre minden $x, y \in K^\times$ esetén
\begin{equation*}
    v(xy) = v(x) + v(y),
\end{equation*}
továbbá, ha $x + y \not = 0$ , akkor

\begin{equation*}
    v(x+y) \geq min\{v(x), v(y)\}.
\end{equation*}
Megállapodás szerint $v(0) = \infty$.

Az első tulajdonság azt mondja, hogy \textit{v} csoporthomomorfizmus $K^\times$ multiplikatív csoportjából $\Gamma$ additív csoportjába. A második tulajdonság az
ultrametrikus háromszög-egyenlőtlenség additív alakja.
\end{definition}

\begin{proposition}
Legyen $v$ értékelés $K$-n. Ekkor:
\begin{enumerate}[label=\arabic*]
    \item $v(1)=0$  és $v(-1) = 0$;
    \item $v(x^{-1})=-v(x)$ \text{minden} $x \in K$ \text{esetén};
    \item ha $v(x)\neq v(y)$, akkor
    $v(x+y)=\min\{v(x),v(y)\}$.
\end{enumerate}
\end{proposition}

\begin{proof}
Az $1 \cdot 1 = 1$ azonosságból $v(1) = v(1) + v(1)$, tehát $v(1) = 0$.
Mivel $(-1)^2 = 1, 2v(-1) = 0$, ezért $v(-1) = 0$ rendezett abel-csoportban,
hiszen abban nincs torzió. Az $xx^{-1} = 1$ azonosságból adódik $v(x)+v(x^{-1}) =
0$, tehát $v(x^{-1}) = -v(x)$.

Tegyük fel, hogy $v(x) < v(y)$. Az \ref{valuation} definícióban szereplő egyenlőtlenség
szerint $v(x + y) \geq v(x)$. Ha szigorú egyenlőtlenség állna fenn, akkor
$v(y) = v((x + y) - x) \geq min\{v(x + y), v(x)\} = v(x)$,
ami ellentmond a kiinduló feltevésnek, következésképpen $v(x+y) = v(x)$.
\end{proof}

\begin{definition}
    Az értékelés értékcsoportja a $v(K^\times) \subseteq \Gamma$ rendezett részcsoport. Az értékelés \textit{triviális}, ha $v(x) = 0$ minden $x \in K^\times$ esetén.
\end{definition}%my part 
\begin{definition}
    Két értékelés $v$ és $v'$ \textit{ekvivalens}, ha létezik $m,n>0$ úgy, hogy $nv = mv'$. 
\end{definition}
\begin{definition}
    Bármely nem triviális értékelés ekvivalens pontosan egy szürjektív értékeléssel. Az ilyen értékeléseket \textit{normalizátor}na nevezzük. 
\end{definition}
% =================================================
\section{Értékelésgyűrűk}

\subsection{Az értékeléshez tartozó gyűrű}

\begin{definition}
Legyen $v$ értékelés a $K$ testen. Az
\[
\Ov = \{x\in K \mid v(x)\ge 0\}
\]
részgyűrűt a $v$ értékelés \textit{értékelésgyűrűjének} nevezzük. Az 
\begin{equation*}
    \m_v = \{x\in K\mid v(x) > 0\}
\end{equation*}
ideált az \textit{értékeléshez tartozó maximális ideálnak} nevezzük. 
\end{definition}
\begin{remark}
    Ha $\Ov$ egy értékelésgyűrűje a $K$ testnek, akkor $K$ bármely részgyűrűje, amely tartalmazza $\Ov$-t szintén értékelésgyűrű. Ez a definícióbol azonnal következik.
\end{remark}

\begin{proposition}
Az $\mathcal O_v$ értékelésgyűrű lokális integritási tartomány, maximális
ideálja $\mathfrak m_v$, egységcsoportja pedig
\[
\mathcal O_v^\times
=
\{x \in K^\times \mid v(x)=0\}.
\]
Továbbá $\operatorname{Frac}(\mathcal O_v)=K$.
\end{proposition}

\begin{proof}
Az értékelés tulajdonságaiból közvetlenül következik, hogy $\mathcal O_v$
zárt az összeadásra és szorzásra, továbbá $1\in\mathcal O_v$. Mivel
$\mathcal O_v \subseteq K$, ezért $\mathcal O_v$ is integritási tartomány.

Ha $x\in\mathcal O_v$ és $v(x)=0$, akkor $v(x^{-1})=0$, tehát
$x^{-1}\in\mathcal O_v$, vagyis $x$ egység. Ha $v(x)>0$, akkor
$v(x^{-1})<0$, tehát $x^{-1}\notin\mathcal O_v$, azaz $x$ nem egység.
Így az egységek pontosan a $0$ értékelésű elemek, a nemegységek pedig
pontosan $\mathfrak m_v$ elemei. Egy gyűrű pontosan akkor lokális,
ha a nemegységek halmaza ideál; itt ez az $\mathfrak m_v$ ideálra teljesül.

Végül ha $x\in K^\times$, akkor vagy $v(x)\ge 0$, és ekkor
$x\in\mathcal O_v$, vagy $v(x)<0$, és ekkor $v(x^{-1})>0$, tehát
$x^{-1}\in\mathcal O_v$. Így minden $K$-beli elem két
$\mathcal O_v$-beli elem hányadosaként írható, tehát
$\operatorname{Frac}(\mathcal O_v)=K$.
\end{proof}

\begin{theorem}\label{tetel23}[Értékelésgyűrűk intrinzikus jellemzése]
Legyen $K$ test, és legyen $R\subseteq K$ olyan részgyűrű, amelyre
$\operatorname{Frac}(R)=K$. Ekkor $R$ pontosan akkor valamely
értékelés értékelésgyűrűje, ha minden $x\in K^\times$ esetén
\[
x\in R
\qquad \text{vagy} \qquad
x^{-1}\in R.
\]
\end{theorem}

\begin{proof}
Ha $R=\mathcal O_v$, akkor az állítás az előző bizonyítás végén már
megjelent.

Fordítva tegyük fel, hogy $R$ rendelkezik a fenti tulajdonsággal.
Tekintsük a $K^\times/R^\times$ faktorcsoportot. Egyelőre multiplikatív
jelölést használunk, de rendezett abel-csoportként később additív
jelölésre térhetünk át. A rendezést definiáljuk úgy, hogy
\[
xR^\times \ge yR^\times
\iff
xy^{-1}\in R.
\]

Ez jól definiált, mert egységgel való szorzás nem változtatja meg az
$R$-beli tagságot. A rendezés teljes, hiszen minden $xy^{-1}\in K^\times$
elemre vagy $xy^{-1}\in R$, vagy $(xy^{-1})^{-1}=yx^{-1}\in R$.
Antiszimmetrikus is: ha $xy^{-1}\in R$ és $yx^{-1}\in R$, akkor
$xy^{-1}$ egység $R$-ben, tehát $xR^\times=yR^\times$.
A tranzitivitás és a csoportművelettel való kompatibilitás közvetlenül
következik abból, hogy $R$ gyűrű.

A
\[
v : K^\times \longrightarrow K^\times / R^\times,
\qquad
x \longmapsto xR^\times
\]
leképezés csoporthomomorfizmus. Ha $x+y\neq 0$ és például
$v(x)\ge v(y)$, akkor $xy^{-1}\in R$, ezért
\[
(x+y)y^{-1}=xy^{-1}+1\in R.
\]

Ez azt jelenti, hogy
\[
v(x+y)\ge v(y)=\min\{v(x),v(y)\}.
\]
A másik eset ugyanilyen. Így $v$ értékelés, és definíció szerint
\[
\mathcal O_v
=
\{x\in K \mid v(x)\ge v(1)\}
=
R.
\]
\end{proof}

\begin{remark}
A fenti tétel miatt az értékelésgyűrűt gyakran azzal a
tulajdonsággal is definiálják, hogy minden $x\in K^\times$ elemre
$x$ vagy $x^{-1}$ benne van a gyűrűben. Ez a megfogalmazás nagyon
hasznos, amikor nem egy konkrét képlettel adott értékelésből
indulunk ki, hanem egy lokális gyűrű tulajdonságait vizsgáljuk. Lásd a \ref{valueex} példák esetén is.
\end{remark}

\subsection{Alapvető tulajdonságok}

\begin{proposition}
Egy értékelésgyűrű ideáljai lineárisan rendezettek a tartalmazásra
nézve: ha $I,J\subseteq \mathcal O_v$ ideálok, akkor
$I\subseteq J$ vagy $J\subseteq I$.
\end{proposition}

\begin{proof}
Tegyük fel, hogy $I\not\subseteq J$. Válasszunk
$a\in I\setminus J$ elemet. Megmutatjuk, hogy $J\subseteq I$.
Legyen $b\in J$. Ha $b=0$, nincs mit bizonyítani.
Ha $b\neq 0$, akkor $\frac{a}{b}\notin \mathcal O_v$, mert ellenkező
esetben
\[
a=(\frac{a}{b})b\in J
\]
lenne. Értékelésgyűrűről lévén szó, ebből
$\frac{b}{a}\in \mathcal O_v$ következik. Így
\[
b=(\frac{b}{a})a\in I,
\]
tehát $J\subseteq I$.
\end{proof}

\begin{remark}
    Az előző állítás megfordítása is igaz. Vagyis, ha $\Ov$ egy integritási tartomány, aminek hányadosteste $K$ és az idáljai rendezettek a benfoglalásra nézve, akkor $\Ov$ értékelésgyűrűje lesz a $K$-nak.
\end{remark}
\begin{proof}
    Tegyük fel, hogy az ideálok teljesen rendezettek az $\Ov$ integritási tartományban. Ha $\alpha \in K^\times$ az $\Ov$ hányadostestéből való, akkor $\alpha = \frac{a}{b}$, ahol $a, b \in \Ov^\times$. A feltételezés szerint vagy $(a) \subset (b)$, vagy $(b) \subset (a)$. Így $a = cb$, ahol $c \in \Ov$, vagy $b = da$, ahol $d \in \Ov $. Az első esetben $\frac{a}{b} = c \in \Ov$, a második esetben pedig $\frac{b}{a} = d \in \Ov$. Tehát $\Ov$ egy értékelésgyűrű. 
\end{proof}
\begin{corollary}
    Legyen $\Ov$ a $K$ test egy részgyűrűje. $\Ov$ akkor és csak akkor értékelésgyűrű, ha
$a, b\in \Ov$ esetén vagy $a$ osztja $b$-t, vagy $b$ osztja $a$-t.
\end{corollary}
\begin{corollary}
    Egy értékelésgyűrűben minden végesen generált ideál principiális.
\end{corollary}
\begin{proof}
     Legyen $I$ egy ideál az $\Ov$ értékelésgyűrűben, és tegyük fel, hogy $I = (a_1, a_2, \dots, a_n)$. Ekkor
vagy $a_1$ osztja $a_2$-t, vagy $a_2$ osztja $a_1$-et. Az általánosság elvesztése nélkül feltételezhetjük, hogy $a_1$ osztja $a_2$-t. Ebben az esetben
$I = (a_1, a_3, …, a_n)$. Ezt a gondolatmenetet folytatva végül azt kapjuk, hogy $I = (a_i)$, valamilyen $i$ esetén. 
\end{proof}

\begin{corollary}
Egy értékelésgyűrű prímideáljai lineárisan rendezettek.
Konkrétan egy értékelésgyűrű spektruma lánc.
\end{corollary}

\begin{proposition}
Minden értékelésgyűrű egészzárt a hányadostestében.
\end{proposition}

\begin{proof}
Legyen $R=\mathcal O_v$, $K=\operatorname{Frac}(R)$, és tegyük fel,
hogy $x\in K$ egész $R$ fölött. Ha $x\notin R$, akkor
$x^{-1}\in R$ és $v(x)<0$. Mivel $x$ egész $R$ fölött,
teljesít egy olyan egyenletet, hogy
\[
x^n+a_{n-1}x^{n-1}+\cdots+a_1x+a_0=0,
\qquad
a_i\in R.
\]

Osszunk $x^{n-1}$-gyel:
\[
x
=
-(a_{n-1}+a_{n-2}x^{-1}+\cdots+a_0x^{-(n-1)}).
\]

A jobb oldal $R$-ben van, hiszen $a_i\in R$ és
$x^{-1}\in R$. Ez ellentmond $x\notin R$-nek.
Tehát $x\in R$, vagyis $R$ egészzárt.
\end{proof}

%my work
\subsection{Példák értékelésgyűrűkre} \label{valueex}
\begin{enumerate}
    \item Bármely $K$ test. Valóban, hiszen $Frac(F) = F$, és így a \ref{tetel23} tétel alapján kész. Ekkor a $K$-hoz tartozó értékelés a triviális értékelés.  
    \item Az egész számok lokalizációja a $(p)$ prímideálban: $\mathbb{Z}_{(p)} = \{ \frac{a}{b} \mid a,b \in \mathbb{Z}, p \nmid b\}$.
    \\ Legyen $x \in Frac (\mathbb{Z}_{(p)}),\, x = \frac{a}{b}$. Tudunk választani olyan $a,b$ egészeket, hogy $p \mid a$ vagy $p|b$, de $p$ nem osztja egyszerre $a$-t is és $b$-t is. Tehát $x \in \mathbb{Z}_{(p)}$ vagy $x^{-1} \in \mathbb{Z}_{(p)}$ 
    \item $\mathbb{C}$ feletti polinomok lokalizációja az $(x)$ prímideálban: \\ $\mathbb{C}[x]_{(x)} = \{ \frac{f(x)}{g(x)} \mid f,g \in \mathbb{C}[x],\, g(0)  \neq 0\}$.
    \item (Nem példa) $\mathbb{Z}$ nem értékelésgyűrű, hiszen például $\frac{3}{2} = (\frac{2}{3})^{-1} \in \mathbb{Q} = Frac(\mathbb{Z})$, de egyik sem egész szám.
    \item  (Nem példa) A $\mathbb{C}$ feletti kétváltozós polinomok lokalizációja: \\$\mathbb{C}[x,y]_{(x,y)} = \{ \frac{f(x,y)}{g(x,y)} \mid f,g \in \mathbb{C}[x,y],\, g(0,0)  \neq 0\}$ bár lokális gyűrű, mégsem értékelésgyűrű. Például $\frac{x}{y} = (\frac{y}{x})^{-1} \in Frac(\mathbb{C}[x,y]_{(x,y)})$, de a gyűrűben egyik sincs benne.
\end{enumerate}

\end{document}
% =================================================
